\section{Job Design and the Firm's Long Run Problem}
\label{app:jobdesign}

The main text studies the firm's AI deployment problem: choosing how to chain and automate steps so as to minimize the cost of production, $\sum_b \skillcost{b}\,\timecost{b}$, where each task is staffed by a dedicated worker possessing the skill it requires (Section~\ref{sec:model.jobs}).
This appendix develops the firm's long-run problem, in which it additionally bundles tasks into \emph{jobs}.
A single worker may perform several tasks, but must then hold all of their skills throughout, so the cost of a job is its total skill requirement times its total time.
Bundling economizes on the coordination costs incurred when work passes from one worker to another, at the expense of these broader skill requirements. 
The firm chooses the AI deployment and the job design jointly in the long-run problem to trade these forces off.
Figure~\ref{fig:hierarchical_production} illustrates the resulting three-layer structure of steps, tasks, and jobs in the long-run optimization problem.

\begin{figure}[ht!]
  \begin{center}
  \caption{Illustrative Example of a Firm's Production in the Long Run}
  \label{fig:hierarchical_production}

  \vspace{0.4cm}
  % Hierarchical Production Visualization
\begin{adjustbox}{center,trim=0pt 0pt 0pt 0pt,clip,margin=-2cm 0cm}
  \begin{tikzpicture}[scale=0.73, transform shape,
    node distance=0.35cm and 0.91cm,
    every node/.style={rectangle, rounded corners, draw, align=center, minimum width=2cm, minimum height=1cm},
    manual/.style={fill=gray!10},
    automated/.style={fill=green!30},
    augmented/.style={fill=orange!30},
    dashedbox/.style={draw, rectangle, rounded corners, dashed, inner sep=0.25cm, line width=1.25pt},
    humanbox/.style={rectangle, rounded corners, draw, align=center, minimum width=2.25cm, minimum height=1cm, fill=olive!20}, % changed to olive
    aibox/.style={rectangle, rounded corners, draw, align=center, minimum width=2.25cm, minimum height=1cm, fill=blue!20},
    jobbox/.style={rectangle, rounded corners, draw, align=center, minimum width=2.25cm, minimum height=1cm, fill=cyan!20},
    industrybox/.style={rectangle, rounded corners, draw, align=center, minimum width=3cm, minimum height=1cm, fill=teal!60},
    >=latex
  ]

    % Top layer (Steps)
    \node[manual] (S1) {Step 1\\(Manual)};
    \node[automated, right=of S1] (S2) {Step 2\\(Automated)};
    \node[automated, right=of S2] (S3) {Step 3\\(Automated)};
    \node[augmented, right=of S3] (S4) {Step 4\\(Augmented)};
    \node[manual, right=of S4] (S5) {Step 5\\(Manual)};
    \node[augmented, right=of S5] (S6) {Step 6\\(Augmented)};
    \node[manual, right=of S6] (S7) {Step 7\\(Manual)};

    % Arrows between steps (specific arrows as curly lines)
    \draw[->, thick] (S1) -- (S2);
    \draw[->, thick] (S2) -- (S3);
    \draw[->, thick] (S3) -- (S4);
    \draw[->, thick, red, decorate, decoration={snake, amplitude=.5mm, segment length=3.05mm}] (S4) -- (S5);
    \draw[->, thick, red, decorate, decoration={snake, amplitude=.5mm, segment length=3.05mm}] (S5) -- (S6);
    \draw[->, thick] (S6) -- (S7);

    % Dashed boxes (top layer, colors updated)
    \node[dashedbox, draw=olive, fit=(S1)] (DS1) {};
    \node[dashedbox, draw=blue, fit=(S2)(S3)(S4)] (DS2-4) {};
    \node[dashedbox, draw=olive, fit=(S5)] (DS5) {};
    \node[dashedbox, draw=blue, fit=(S6)] (DS6) {};
    \node[dashedbox, draw=olive, fit=(S7)] (DS7) {};

    \path (S1.west) -- (S7.east) coordinate[midway] (CenterTop);

    % Middle layer (Tasks, human task color updated)
    \node[humanbox, below=3.2cm of CenterTop, xshift=-6cm] (T1) {Task 1\\(Human)};
    \node[aibox, right=of T1] (T2) {Task 2\\(AI Chain)};
    \node[humanbox, right=of T2] (T3) {Task 3\\(Human)};
    \node[aibox, right=of T3] (T4) {Task 4\\(AI Chain)};
    \node[humanbox, right=of T4] (T5) {Task 5\\(Human)};

    % Arrows between tasks (specific arrows as curly lines)
    \draw[->, thick] (T1) -- (T2);
    \draw[->, thick, red, decorate, decoration={snake, amplitude=.5mm, segment length=3.05mm}] (T2) -- (T3);
    \draw[->, thick, red, decorate, decoration={snake, amplitude=.5mm, segment length=3.05mm}] (T3) -- (T4);
    \draw[->, thick] (T4) -- (T5);

    % Dashed boxes (middle layer)
    \node[dashedbox, draw=cyan, fit=(T1)(T2)] (DT1-2) {};
    \node[dashedbox, draw=cyan, fit=(T3)] (DT3) {};
    \node[dashedbox, draw=cyan, fit=(T4)(T5)] (DT4-5) {};

    % Job layer nodes
    \node[jobbox, below=2.25cm of T3] (J2) {Job 2};
    \node[jobbox, left=of J2] (J1) {Job 1};
    \node[jobbox, right=of J2] (J3) {Job 3};

    % Arrows between jobs (specific arrows as curly lines)
    \draw[->, thick, red, decorate, decoration={snake, amplitude=.5mm, segment length=3.05mm}] (J1) -- (J2);
    \draw[->, thick, red, decorate, decoration={snake, amplitude=.5mm, segment length=3.05mm}] (J2) -- (J3);

    % Arrows connecting layers (unchanged)
    \draw[->, thick, cyan] (DT1-2.south) -- (J1.north);
    \draw[->, thick, cyan] (DT3.south) -- (J2.north);
    \draw[->, thick, cyan] (DT4-5.south) -- (J3.north);

    \draw[->, thick, olive] (DS1.south) -- (T1.north);
    \draw[->, thick, blue] (DS2-4.south) -- (T2.north);
    \draw[->, thick, olive] (DS5.south) -- (T3.north);
    \draw[->, thick, blue] (DS6.south) -- (T4.north);
    \draw[->, thick, olive] (DS7.south) -- (T5.north);

  \end{tikzpicture}
\end{adjustbox}
  \end{center}

  \vspace{0.4cm}
  \footnotesize{\emph{Notes:}
  This figure illustrates how a firm aggregates steps into tasks and tasks into jobs.
  The top layer shows $m=7$ steps executed in manual (gray), automated (green), or augmented (orange) modes, with dashed boxes indicating task groupings.
  The middle layer aggregates steps into $n=5$ tasks, classified as human-executed tasks (olive) or AI chains (purple), with cyan dashed boxes indicating job groupings.
  The bottom layer aggregates tasks into $p=3$ jobs (cyan), each assigned to a separate worker.
  Vertical colored arrows show aggregation across layers, and horizontal arrows show production flow within each layer.
  Red curly arrows mark hand-off costs incurred at job boundaries, and can be traced from the bottom layer to the corresponding tasks and steps above.
  }
\end{figure}


\subsection{Jobs, Wages, and Hand-off Costs}
\label{sec:app.jobs}

A job is a subsequence of tasks that are assigned to a single human worker.
Each task $T_b$ in the task sequence $\T = \{T_1, \dotsc, T_n\}$ has a skill cost $\skillcost{b}$ and a time cost $\timecost{b}$ that depend on whether the task is completed manually or with the aid of AI (Sections~\ref{sec:model.steps}--\ref{sec:model.tasks}).
A job $J$ is a contiguous subsequence of tasks, say $J = (T_b, \dotsc, T_{b+\ell})$ for some $b$ and $\ell$.
We say that a partition $\mathcal{J} = (J_1, \dotsc, J_p)$ of all tasks into jobs is a \emph{job design}.

The firm hires one worker for each job in its job design $\mathcal{J}$.
The worker for job $J_j$ completes all tasks associated with their job.
The total time needed to complete all tasks in job $J_j$ is $\sum_{T_b \in J_j} \timecost{b}$.

A worker's wage is determined by the total skill required to complete their job.
The total skill required to complete the tasks in job $J_j$ is $\sum_{T_b \in J_j} \skillcost{b}$.
We then assume that a worker's wage is proportional to the skill level of their job.
That is,
\begin{align}
\label{eq:wage_jobdesign}
\text{Wage}_j = \sum_{T_b \in J_j} \skillcost{b}.
\end{align}

The firm must pay workers their wage per unit of time regardless of which tasks they are assigned to complete at any given moment.
While we allow the firm to specialize its job design and partition its workflow into distinct jobs, we acknowledge that there are inherent inefficiencies and frictions that are introduced when splitting work among multiple workers.
Absent such frictions in our model, it would always be optimal for a firm to specialize its workers as much as possible, assigning a distinct worker to each task of its workflow. 
We therefore introduce \emph{hand-off costs} to a job design, as follows.
In addition to the time required completing tasks, a worker may need to spend time handing off their work to another worker if $J_j$ is not the final job in the job design $\mathcal{J}$.
We will assume throughout that the hand-off time of a worker assigned to job $J_j$ depends only on the final step of job $J_j$ (see red curly arrows in Figure \ref{fig:hierarchical_production}).
That is, conditional on where the hand-off occurs in the workflow, its cost does not otherwise depend on previous hand-off events.
Given step $s_i$, we write $\handofftime{i}$ to denote the additional hand-off time spent by a worker for whom the last step of their job is $s_i$.
For notational convenience we define the hand-off time for the final step $s_m$ to be zero: $\handofftime{m} = 0$.
Also for notational convenience, we will write $\handofftime{}(J_j)$ to denote the hand-off time of job $J_j$, which recall is equal to $\handofftime{i}$ if $s_i$ is the final step in job $J_j$.
Thus the total time needed to complete job $J_j$, including hand-off costs, is
\[
\handofftime{}(J_j) + \sum_{T_b \in J_j} \timecost{b}.
\]

Taking into account both time and skill requirements, the total wage bill paid to a worker assigned to job $J_j$, per unit of output, is
\begin{align}
\label{eq:wagebill}
\text{WageBill}_j = \Biggl(\sum_{T_b \in J_j} \skillcost{b} \Biggr) \Biggl(\handofftime{}(J_j) + \sum_{T_b \in J_j} \timecost{b} \Biggr).
\end{align}
Equation~\eqref{eq:wagebill} highlights two opposing forces that shape a job's boundaries.  
Adding more tasks to a worker's job increases the cumulative skill requirement and thus the wage rate for that worker.  
However, if tasks are kept separate, workers must incur additional hand-off costs at each job boundary.  
These two effects, higher wages from combining tasks versus higher hand-off costs from keeping them separate, jointly determine the optimal degree of task aggregation.  
Figure~\ref{fig:intro_task_illustration} illustrates these opposing forces in a minimal two-task workflow.
We explore this trade-off in greater detail in Section~\ref{sec:model_discussion}.

\begin{figure}[ht!]
    \begin{center}
    \caption{Illustration of the Trade-off between Worker Specialization and Coordination in Task Assignment}
    \label{fig:intro_task_illustration}
    \includegraphics[width=1\linewidth]{plots/job_design.png}
    \end{center}
    \footnotesize
    \emph{Notes:} The blue bounded rectangles indicate tasks.
    The height of each rectangle corresponds to the task's skill requirement ($\skillcost{}$) and its width to the time requirement ($\timecost{}$).
    The shaded areas represent the wage bills of jobs.
    In the left panel, tasks are assigned to two specialized workers but a hand-off cost (the pink rectangle) is introduced due to coordination frictions between them.
    In the right panel, tasks are bundled into a single job performed by one worker, eliminating hand-off costs but requiring a worker skilled enough to perform both tasks and compensated with a higher per-unit-time wage rate.
\end{figure}

Table~\ref{tab:jobdesign_parameters} collects the additional notation introduced for job design.
\begin{table}[ht!]
\centering
\caption{Additional Notation for Job Design}
\label{tab:jobdesign_parameters}
\begin{tabularx}{\textwidth}{l X}
\toprule
\textbf{Symbol} & \textbf{Description} \\
\midrule
$J_j$ & Job $j$: a contiguous block of tasks assigned to a single worker \\
$\mathcal{J} = (J_1, \dotsc, J_p)$ & Job design; a partition of the tasks into $p$ jobs \\
\addlinespace
$\text{Wage}_j$ & Wage rate of job $J_j$'s worker, equal to the total skill it requires, $\sum_{T_b \in J_j} \skillcost{b}$ \\
$\handofftime{i}$ & Hand-off time incurred when a worker's final step is $s_i$ \\
$\handofftime{}(J_j)$ & Hand-off time of job $J_j$, equal to $\handofftime{i}$ for its final step $s_i$ \\
$\text{WageBill}_j$ & Per-unit-output labor cost of job $J_j$: $\bigl(\sum_{T_b \in J_j} \skillcost{b}\bigr)\bigl(\handofftime{}(J_j) + \sum_{T_b \in J_j} \timecost{b}\bigr)$ \\
\bottomrule
\end{tabularx}
\end{table}


\subsection{Firm's Long-Run Optimization Problem}
\label{app:jobdesign.longrun}

In the long-run optimization problem, the firm anticipates the adjustment of worker wages to fit the skill requirements of each job and sets job responsibilities accordingly.
This contrasts with the short- and medium-run AI strategy deployment problems of Section~\ref{sec:model.jobs}, which minimize the skill-weighted cost of production while fixing the assignment of tasks to workers. 
The long-run problem additionally optimizes how tasks are bundled into jobs.

Given the sequence of steps $\mathcal{S} = \{s_1, \dotsc, s_m\}$, the firm can design both the partition $\T$ of steps into tasks and the partition $\mathcal{J}$ of tasks into jobs.
The choice of $\T$ determines the time and skill requirements of each task.
Given $\T$, the choice of $\mathcal{J}$ then determines worker wage rates and time needed per unit of output (including hand-off costs).
Formally, we can write $\mathcal{P}(X)$ for the set of partitions of a sequence $X$ into contiguous subsequences.  Then the long-run optimization problem faced by the firm can be expressed as
\begin{align}
\label{eq:totalcost_with_handoff}
\min_{\T \in \mathcal{P}(\mathcal{S})} \
\min_{\mathcal{J} \in \mathcal{P}(\T)} \ 
\text{TotalCost}(\mathcal{J}; \T) 
= \sum_{J_j \in \mathcal{J}} \text{WageBill}_j = 
\sum_{J_j \in \mathcal{J}} \Biggl[\Biggl(\sum_{T_b \in J_j} \skillcost{b} \Biggr) \Biggl(\handofftime{}(J_j) + \sum_{T_b \in J_j} \timecost{b} \Biggr)\Biggr]
\end{align}
where recall that, for each $T_b \in \T$, $\timecost{b}$ and $\skillcost{b}$ are as described in Section~\ref{sec:model.tasks} and can depend on the selected mode of operation for individual steps.

This problem has a similar flavor to the medium-run problem of Proposition~\ref{prop:mediumrun_dp}, with one additional layer: as steps are assigned to tasks, the resulting tasks must also be assigned to workers along the way by considering the hand-off costs incurred at the boundary of jobs.
Proposition~\ref{prop:totalcost_optimization_dp} shows that, like the medium-run problem, the long-run problem admits a polynomial-time approximation scheme.

\begin{proposition}[Long-run AI Deployment and Job Design]
\label{prop:totalcost_optimization_dp}
Fix a workflow with $m$ steps whose skill, time, and hand-off costs lie in $[1/B, B]$ for some $B > 0$.
For any $\epsilon > 0$, a pair of AI strategy $\T$ and job design $\J$ whose long-run cost~\eqref{eq:totalcost_with_handoff} is within a factor $(1+\epsilon)$ of the minimum can be computed in time $O(m^2 \epsilon^{-2} \log^2(mB))$ via dynamic programming.
\end{proposition}

\begin{proof}
We prove the proposition following the same approach as the proof of the medium-run problem (Proposition~\ref{prop:mediumrun_dp}): we first state the cost minimization problem as a recursive problem and then use dynamic programming to approximate the optimal solution.
Since much of the construction carries over, we focus on what is new, the assignment of tasks to workers and the associated hand-off costs, and refer to the earlier proof for details that are unchanged.

\emph{Recursive formulation.}
Consider a workflow with $m$ steps $\mathcal{S} = \{s_1, s_2, \ldots, s_m\}$.
For $0 \leq i \leq m$, let $V(i)$ denote the minimum cost required to complete steps $1$ through $i$ using one or more jobs, including any hand-off costs to a subsequent worker.
Note then that $V(m)$ is the solution to the desired optimization problem, recalling that the hand-off time for the final job will always be $0$.
Note also that it is implicitly assumed in the definition of $V(i)$ that it optimizes over job designs for steps $1$ through $i$ in which the final job terminates after the completion of step $i$.

For job designs whose final job does not necessarily terminate after step $i$ but rather continues onward to subsequent steps, we extend the auxiliary function $R(i,\; \skillcost{},\; \timecost{})$ from the proof of Proposition~\ref{prop:mediumrun_dp}.
As before, $R(i,\; \skillcost{},\; \timecost{})$ denotes the minimum cost of completing steps $1$ through $i$ when the worker currently accumulating work, whom we call the ``active worker,'' has already been assigned tasks with total skill cost $\skillcost{}$ and total time cost $\timecost{}$, where $\timecost{}$ now also includes any hand-off cost for the active worker.
What is new is that the minimum is taken over job designs as well as AI deployment strategies: the recursion may either expand the active worker's job or close it and assign the remaining steps to other workers.
The base case $R(0,\; \skillcost{},\; \timecost{}) = \skillcost{}\,\timecost{}$ is unchanged, and as before we can think of step $i$ as the highest-indexed step that has not yet been assigned.

The two functions are linked, for all $0 \leq i \leq m$, by
\[ V(i) = R(i,\; 0,\; \handofftime{i}). \]
This is because, in $V(i)$, all job designs must end with a job that completes at step $i$ and hands off to the following worker.
Incurring this hand-off cost is equivalent to assigning the active worker an additional task with zero skill requirement and a time cost of $\handofftime{i}$ right after step $i$ itself, which is exactly the scenario $R(i,\; 0,\; \handofftime{i})$ represents.

For each $i \geq 1$, the minimum cost function $R(i,\; \skillcost{},\; \timecost{})$ satisfies the recursive relation
\begin{align*}
R(i,\; \skillcost{},\; \timecost{}) = \min\Biggl\{
& \ \underbrace{\skillcost{}\,\timecost{} + V(i)}_{\text{job closed}}, \ \underbrace{R\bigl(i-1,\; \skillcost{} + \manualSkill{i},\; \timecost{} + \manualTime{i}\bigr)}_{\text{Step $i$ is manual}}, \ \underbrace{\min_{0 \le r < i}\ R\!\left(r,\; \skillcost{} + \AIskill{i},\; \timecost{} + \frac{\AItime{i}}{\prod_{s=r+1}^{i} q_s}\right)}_{\text{Step $i$ optimally chained}}
\Biggr\}.
\end{align*}
In words, when assigning step $i$, the three options the firm considers are:
\begin{itemize}
    \item \textbf{Job closure.}
    The term $\skillcost{}\,\timecost{} \; + \; V(i)$ corresponds to not adding any further steps to the job of the active worker.
    All steps $1$ through $i$ will be completed by other workers.
    The active worker's total cost is then $\skillcost{}\,\timecost{}$, and $V(i)$ is the total cost of completing steps $1$ through $i$ (including hand-off costs to the active worker).
    This option is the margin new to the long run.

    \item \textbf{Manual step.}
    The term $R(i-1,\; \skillcost{} + \manualSkill{i},\; \timecost{} + \manualTime{i})$ corresponds to designating step $i$ for manual completion, and adding the corresponding (singleton) task.
    As in the medium run, the manual execution costs $(\manualSkill{i}, \manualTime{i})$ of step $i$ are added to the accumulated costs of tasks already assigned to the active worker, extending their job to include step $i$ as well.

    \item \textbf{AI chain.}
    The term $\min_{0 \le r < i} \ R\left(r,\; \skillcost{} + \AIskill{i},\; \timecost{} + \frac{\AItime{i}}{\prod_{s=r+1}^{i} q_s}\right)$ corresponds to creating an AI chain task for steps $r+1$ through $i$, with step $i$ being augmented and steps $r+1$ through $i-1$ automated, and assigning the resulting task to the job of the active worker.
    Step $i$ is thus chained with zero or more preceding (automated) steps, and the associated AI management costs are added to the active worker's skill and time costs.
    The cut $r$ is the highest-indexed step that is not added to this AI chain, which can be any value between $0$ and $i-1$.
\end{itemize}
Note that in evaluating the AI chain option, step $i$ can only be augmented (rather than automated), as each step of the recursion adds a full AI chain to the active worker's job; AI strategies in which step $i$ is automated are considered when performing the calculations $R(r,\; \skillcost{},\; \timecost{})$ with $r > i$.
The cost of the optimal joint AI strategy and job design for the full sequence of $m$ steps is then $V(m) = R(m,\; 0,\; 0)$.

\emph{Discretization and Computation Time.}
As in the proof of Proposition~\ref{prop:mediumrun_dp}, we fill in the table of values $R(i,\; \skillcost{},\; \timecost{})$ by dynamic programming after discretizing the skill and time coordinates to powers of $(1+\epsilon)$; only the range of the discretization requires a new argument, to account for hand-off costs.
Completing every step manually with a separate worker incurs, for each step, a time and hand-off cost of at most $B$ each and a skill cost of at most $B$, for a total cost of at most $2B^2$ per step, so no solution containing a job with total cost above $2mB^2$ need be considered.
Since skill costs combine additively and every non-zero skill or time cost is at least $1/B$, it suffices to tabulate skill costs in $[1/B,\, mB]$ and time costs in $[1/B,\, 2mB^3]$.
Each range contains $O(\epsilon^{-1}\log(mB))$ powers of $(1+\epsilon)$, so the table has $O(m\,\epsilon^{-2}\log^2(mB))$ entries, each filled in time $O(m)$ via the recursion above, for a total runtime of $O(m^2\,\epsilon^{-2}\log^2(mB))$; recording which option attains each minimum lets the firm read the optimal AI strategy and job design off $V(m) = R(m,\; 0,\; 0)$.
The discretization error is bounded exactly as before: skill and time costs accumulate additively along the manual and chaining options, so rounding down to powers of $(1+\epsilon)$ maintains a multiplicative error of at most $(1+\epsilon)$ in each coordinate, while each job's wage bill is formed by the single multiplication $\skillcost{}\,\timecost{}$ when the first option closes it, contributing a factor of at most $(1+\epsilon)^2$.
Rescaling $\epsilon$ by a constant therefore yields the claimed $(1+\epsilon)$ approximation.
\end{proof}

\subsection{Hand-off Costs and the Limits of Worker Specialization}
\label{sec:model_discussion}

Here we discuss how hand-off costs determine the optimal organization of production and limit full worker specialization in the firm's long-run problem \eqref{eq:totalcost_with_handoff}.
We leverage insights from a numerical example as well as a geometric illustration, which together show how task aggregation balances higher wages against lower coordination costs.

Consider a workflow that, under a fixed AI deployment strategy $\T$, has $n=3$ tasks with the following cost parameters:
\[
\begin{array}{c|ccc}
\textbf{Task} & \skillcost{b} & \timecost{b} & \handofftime{b} \\ \hline
1 & 3  & 1   & 3 \\
2 & 1  & 2   & 0.5 \\
3 & 2  & 2   & 0 \\
\end{array}
\]

With three tasks, there are four ways the firm can design jobs.\footnote{
More generally, a workflow with $n$ tasks has $2^{(n-1)}$ unique job designs.
} 
Denoting jobs by square brackets, with three tasks the four possible job designs are 
$\{[1][2][3], \quad [1,2][3], \quad [1][2,3], \quad [1,2,3]\}$.
The cost of these job designs for the example above is given in Table~\ref{tab:job_design}.
\begin{table}[ht!]
    \caption{Role of Hand-off Costs in Organizational Structure}
    \vspace{-0.4cm}
    \begin{center}
    \resizebox{0.9\textwidth}{!}{%
    \begin{tabular}{c}
    \begin{minipage}{\textwidth}
    \centering
    Panel (a): Without Hand-off Costs \\[1ex]
    \begin{tabular}{|
  >{\centering\arraybackslash}m{1.7cm}|
  >{\centering\arraybackslash}m{0.9cm}|
  >{\centering\arraybackslash}m{1.8cm}|
  >{\centering\arraybackslash}m{1.8cm}|
  >{\centering\arraybackslash}m{1.8cm}|
  >{\centering\arraybackslash}m{1.6cm}|
  >{\centering\arraybackslash}m{1.7cm}|
  >{\centering\arraybackslash}m{1.7cm}|
  }
\hline
\textbf{Job Design} & \textbf{Job} & \textbf{Job Tasks} & $\sum\hccost{b}$ & $\sum\timecost{b}$ & \textbf{Job Cost} & \textbf{Total Cost} & \textbf{Optimal Design} \\
\hline
\multirow{3}{*}{[1][2][3]} 
  & 1 & $\{1\}$   & 3 & 1 & 3 & \multirow{3}{*}{9} & \multirow{3}{*}{\checkmark} \\
\cline{2-6}
  & 2 & $\{2\}$   & 1 & 2 & 2 &  &  \\
\cline{2-6}
  & 3 & $\{3\}$   & 2 & 2 & 4 &  &  \\
\hline
\multirow{2}{*}{[1,2][3]}
  & 1 & $\{1,2\}$ & 4 & 3 & 12 & \multirow{2}{*}{16} &  \\
\cline{2-6}
  & 2 & $\{3\}$   & 2 & 2 & 4 &  &  \\
\hline
\multirow{2}{*}{[1][2,3]}
  & 1 & $\{1\}$   & 3 & 1 & 3 & \multirow{2}{*}{15} &  \\
\cline{2-6}
  & 2 & $\{2,3\}$ & 3 & 4 & 12 &  &  \\
\hline
[1,2,3] & 1 & $\{1,2,3\}$ & 6 & 5 & 30 & 30 &  \\
\hline
\end{tabular}
    \end{minipage} \\[21ex]
    
    \begin{minipage}{\textwidth}
    \centering
    Panel (b): Including Hand-off Costs \\[1ex]
    \begin{tabular}{|
  >{\centering\arraybackslash}m{1.7cm}|
  >{\centering\arraybackslash}m{0.9cm}|
  >{\centering\arraybackslash}m{1.8cm}|
  >{\centering\arraybackslash}m{1.8cm}|
  >{\centering\arraybackslash}m{1.8cm}|
  >{\centering\arraybackslash}m{1.6cm}|
  >{\centering\arraybackslash}m{1.7cm}|
  >{\centering\arraybackslash}m{1.7cm}|
  }
\hline
\textbf{Job Design} & \textbf{Job} & \textbf{Job Tasks} & $\sum\hccost{b}$ & $\handofftime{} + \sum\timecost{b}$ & \textbf{Job Cost} & \textbf{Total Cost} & \textbf{Optimal Design} \\
\hline
\multirow{3}{*}{[1][2][3]} 
  & 1 & $\{1\}$   & 3 & 4 & 12 & \multirow{3}{*}{18.5} &  \\
\cline{2-6}
  & 2 & $\{2\}$   & 1 & 2.5 & 2.5 &  &  \\
\cline{2-6}
  & 3 & $\{3\}$   & 2 & 2 & 4 &  &  \\
\hline
\multirow{2}{*}{[1,2][3]}
  & 1 & $\{1,2\}$ & 4 & 3.5 & 14 & \multirow{2}{*}{18} & \multirow{2}{*}{\checkmark} \\
\cline{2-6}
  & 2 & $\{3\}$   & 2 & 2 & 4 &  &  \\
\hline
\multirow{2}{*}{[1][2,3]}
  & 1 & $\{1\}$   & 3 & 4 & 12 & \multirow{2}{*}{24} &  \\
\cline{2-6}
  & 2 & $\{2,3\}$ & 3 & 4 & 12 &  &  \\
\hline
[1,2,3] & 1 & $\{1,2,3\}$ & 6 & 5 & 30 & 30 &  \\
\hline
\end{tabular}
    \end{minipage}
    \end{tabular}
    }
    \label{tab:job_design}
    \end{center}
    \footnotesize{\emph{Notes:} This table reports total production costs for each possible job design in an example workflow with a fixed automation strategy and $n = 3$ tasks. 
    The cost parameters of the tasks are given by $(\skillcost{b},\timecost{b},\handofftime{b})=\{(3,1,3),(1,2,0.5),(2,2,0)\}$.
    In the absence of hand-off frictions, full specialization, corresponding to job design [1][2][3], is cost-minimizing. 
    Once hand-off costs are taken into account, the optimal organizational structure changes to job design [1,2][3], which avoids the large coordination cost between tasks 1 and 2 at the expense of employing a higher-cost (more skilled) worker to complete those tasks.
    }
\end{table}
In the absence of hand-off costs in Panel (a), the optimal job design corresponds to full specialization: each task forms a separate job assigned to a different worker.
When hand-off costs are introduced in Panel (b), however, the optimal design changes to one in which tasks 1 and 2 are combined into a single job, while task 3 remains separate.
This structure is optimal because it avoids the large hand-off cost that would otherwise arise between tasks 1 and 2 due to the large hand-off time value $\handofftime{1}=3$.
By aggregating tasks 1 and 2, the firm pays a higher wage to a more skilled and more versatile worker but eliminates the costly coordination between the two tasks, lowering total production cost.
Equation~\eqref{eq:wagebill} captures this trade-off explicitly: 
adding more tasks to a worker's job raises the wage rate through the skill term $\sum_{T_b \in J_j} \skillcost{b}$, 
while keeping tasks separate raises total cost through additional hand-offs $\handofftime{}(J_j)$.
These two effects jointly determine the optimal degree of task aggregation.

The cost of different job designs admits a geometric interpretation as well.  
Figure~\ref{fig:job_design} visualizes the workflow of the example in Table~\ref{tab:job_design} as stacked rectangles, where each rectangle's height corresponds to $\skillcost{b}$ and its width to $\timecost{b}$.
\begin{figure}[ht!]
\begin{center}
\caption{Geometric Interpretation of Job Designs} \label{fig:job_design}
\begin{subfigure}[b]{0.49\textwidth}
    \caption{Job Design Without Hand-off Costs}
    \includegraphics[width=\textwidth]{plots/combined_grid_no_handoff.png}
\end{subfigure}
\begin{subfigure}[b]{0.49\textwidth}
    \caption{Job Design With Hand-off Costs}
    \includegraphics[width=\textwidth]{plots/combined_grid_with_handoff.png}
\end{subfigure}
\end{center}
\footnotesize{\emph{Notes:} Each panel visualizes the workflow as stacked rectangles, where the area of each job's bounding box represents its wage bill.  
Panel (a) sets hand-off costs to zero, in which case it is optimal to assign one task per job.
Panel (b) introduces hand-off costs, which alters the cost structure and makes the [1,2][3] design cost-minimizing.
}
\end{figure}
Panel~(a) depicts the job design when hand-off costs are zero, Panel~(b) shows the same workflow with hand-off costs.  
Visually, optimal job design balances two opposing forces.  
Combining tasks into a single job eliminates the pink rectangles representing additional hand-off costs between tasks, but creates inefficiencies due to the higher skill level of the worker required to do the job. 
Conversely, splitting tasks into separate jobs eliminates the intra-job inefficiencies but incurs hand-off time costs at each job boundary.
These boundary costs scale with the cumulative human capital of the switching worker, so the height of the pink hand-off rectangle in Panel (b) equals the sum of human capital heights for that job. 
Cost-effective job design must balance these effects.

We have modeled hand-off costs as fixed and not influenced by the introduction of AI. 
In reality, AI may reduce the cost of hand-offs between workers as well. 
This might occur because the communication burden of handing off tasks could be partially automated through AI assistance. 
While we do not model this explicitly, we expect that such a force would reduce the impact of hand-off costs and result in a higher degree of worker specialization.\footnote{
A useful interpretation of hand-off costs is that they consist of two components: (i) peripheral coordination work that can be facilitated by AI (e.g., drafting messages, summarizing prior work, formatting documentation), and (ii) a purely manual remainder that relies on tacit, interpersonal, or social knowledge and is not easily executable by AI \citep{deming2017growing}. 
For tractability, we abstract from the first component in our framework and treat its benefits as absorbed into execution and AI management time costs. 
The hand-off cost in the model should therefore be read as capturing the second, irreducibly human component, even though AI-facilitated reductions in the peripheral component may matter in practice.
}

\begin{comment}
\subsection{Impact of AI Deployment on Worker Skill and Specialization}
\label{sec:discussion.skills}

Because the firm minimizes the skill-weighted cost of production, AI lowers cost by reducing the time required for a task, the skill required to perform it, or both.
The skill channel can dampen or even reverse the intuition that gains from AI derive primarily from faster execution: AI may be worth deploying on a step even when it takes longer, provided it sufficiently lowers the skill required, as the following example shows.

\begin{example}
    Consider a job consisting of a single step, with the parameters below:
\begin{center}
\begin{tabular}{c|ccc}
\toprule
\makecell{Step\\$(i)$} & \makecell{Manual cost\\$(\manualSkill{i}, \manualTime{i})$} & \makecell{AI management cost\\$(\AIskill{i}, \AItime{i})$} & \makecell{AI success probability\\$(q_i)$} \\
\midrule
1 & $(5, 5)$ & $(1, 1)$ & $1/8$ \\
\bottomrule
\end{tabular}
\end{center}
    Despite the task taking longer to complete with AI (due to the low likelihood of success on any individual attempt), the reduced skill requirement of managing the AI means that it is firm-optimal to employ AI augmentation.
\end{example}

The impact of AI on worker skill reflects two common narratives about the impact of AI on work: first, that AI can automate mundane or repetitive tasks and allow high-skill workers to concentrate more of their time on meaningful high-skill tasks; second, that AI could lead to deskilling as high-skilled labor for manual task completion is replaced with low-skill AI management. 
Our framework unifies these two perspectives by endogenizing (a) the nature (i.e., time and skill requirements) of work to be automated and (b) how the AI strategy impacts requisite worker skill.  
If replacing a given sequence of production steps in given job with an AI chain is beneficial, in the long term, then either the total time needed to complete the steps is shortened, the total skill needed to manage the AI automation is reduced relative to completing the steps manually, or both.  
For an example of the latter, a step with low skill requirement but high time requirement (such as Step 2 in Figure~\ref{fig:job_design}) may be optimally augmented by AI even when doing so increases the skill needed for AI management, if the time savings are substantial enough.  
In this case, the skill required to complete the same job might increase: AI augmentation can be viewed as complementing worker capabilities.  
Alternatively, it may also be beneficial to employ AI on production steps with high skill requirements (such as Step 1 in Figure~\ref{fig:job_design}) even if this results in an AI-assisted task that takes longer to complete, as long as it requires substantially less worker skill to manage the AI.  
This can ultimately lead to a deskilling of the labor force assigned to the job as AI capital substitutes for high-skilled human work.  

This discussion has so far focused on impacts within a single job. 
Our framework likewise captures how AI deployment strategy shapes job design and, in particular, patterns of worker specialization across tasks. 
Consider again the example from Figure~\ref{fig:job_design}, which highlights how hand-off costs and the structure of what we refer to as \emph{tent-pole tasks} can influence how production steps are bundled into specialized jobs. 
A tent-pole task is a short but high-skill task that sits immediately next to time-consuming, low-skill ones.\footnote{
This juxtaposition creates a classic friction: if a single highly skilled worker performs all adjacent tasks, a substantial share of their time is spent on low-skill work, whereas assigning the surrounding low-skill tasks to different workers raises hand-off costs when work passes between them. 
Historically, this tension has been an important driver of the division of labor. 
As AI automation alters the time and skill profiles of individual tasks in the production sequence, it changes exactly this trade-off, shifting when it is efficient to have a single worker perform a cluster of tasks versus when specialization across workers remains optimal.
}
The second task in Figure~\ref{fig:tent_pole} is an example of a tent-pole task.
\begin{figure}[ht!]
 \caption{Illustration of Tent‐Pole Tasks} 
 \label{fig:tent_pole}
 \begin{center}
   \includegraphics[width=0.4\textwidth]{plots/tent_pole.png}
 \end{center}
 \begin{footnotesize}
   \emph{Notes:} This figure illustrates tent‐pole tasks.  
   The width and height of each rectangle specify the time ($\timecost{}$) and human capital ($\skillcost{}$) requirements of the task, respectively.
   Tasks 1 and 3 have a high time cost and low human capital requirement (wide and short) whereas Task 2, which is the tent-pole task, exhibits a low time cost and high human capital requirement (narrow and tall).
   The hand-off time costs (\handofftime{}) are assumed to be zero for simplicity.
 \end{footnotesize}
\end{figure}

As AI deployment strategy influences the time and skill profile of tasks, the incentive to specialize workers likewise changes.
To predict the direction of this effect, one hypothesis is that AI deployment will tend to have a normalizing effect on tasks, reverting both skill requirements and time requirements toward the mean and making tasks more ``square-like'' (in terms of our geometric interpretation of job design costs).  
If so, and if hand-off time costs remain unaffected, then (roughly speaking) tent-pole inefficiencies due to grouping tasks together into the same job would be reduced.  
This suggests that AI deployment may reduce worker specialization.

\begin{example}
   Consider a two-step workflow with the parameters below, where $\handofftime{i}$ denotes the hand-off cost incurred if a job boundary is placed after step $i$:
\begin{center}
\begin{tabular}{c|cccc}
\toprule
\makecell{Step\\$(i)$} & \makecell{Manual cost\\$(\manualSkill{i}, \manualTime{i})$} & \makecell{AI management cost\\$(\AIskill{i}, \AItime{i})$} & \makecell{AI success\\probability $(q_i)$} & \makecell{Hand-off cost\\$(\handofftime{i})$} \\
\midrule
1 & $(2, 4)$ & $(2, 2)$ & $1$ & $5$ \\
2 & $(4, 2)$ & $(2, 2)$ & $1$ & -- \\
\bottomrule
\end{tabular}
\end{center}
   The first step is low-skill but time-intensive and the second step is high-skill but can be completed quickly.
   Combining both steps into a single job (at a labor cost of $(2+4)(4+2) = 36$) is strictly worse than separating into two specialized jobs with a hand-off (at a total cost of $(2)(4+5) + (4)(2) = 26$).
   However, augmenting each step separately via AI yields effective skill and time costs of $(2,2)$ for each, so the optimal design combines both augmented steps into a single job at a total cost of $(2+2)(2+2) = 16$, which is better than separating into two distinct jobs (incurring cost $(2)(2+5) + (2)(2) = 18$).
\end{example}

It is also technically possible in our model for AI to increase the amount of specialization in an optimal job design.  
This could happen if AI-augmented steps have higher skill requirements than the corresponding manual steps, leading to an increased need for high-skill workers.  
Even if AI management is assumed to require no more skill than manual completion, an increased deployment of AI can lead to more responsibilities being combined into a single job, leading to a higher worker skill requirement, as the following example shows. 

\begin{example}
   Consider a different two-step workflow with the parameters below, where $\handofftime{i}$ denotes the hand-off cost incurred if a job boundary is placed after step $i$:
\begin{center}
\begin{tabular}{c|cccc}
\toprule
\makecell{Step\\$(i)$} & \makecell{Manual cost\\$(\manualSkill{i}, \manualTime{i})$} & \makecell{AI management cost\\$(\AIskill{i}, \AItime{i})$} & \makecell{AI success\\probability $(q_i)$} & \makecell{Hand-off cost\\$(\handofftime{i})$} \\
\midrule
1 & $(3, 3)$ & $(2, 1/4)$ & $1/8$ & $5$ \\
2 & $(3, 3)$ & $(2, 1/4)$ & $1/8$ & -- \\
\bottomrule
\end{tabular}
\end{center}
   The optimal job design separates these into two separate jobs at a total cost of $(3)(3+5) + (3)(3) = 33$, with each job requiring a worker of skill $3$.
   Under AI augmentation each step has an expected execution time of $(1/4) \times (1/8)^{-1} = 2$, giving effective skill and time costs of $(2,2)$.
   In this case, the optimal design employs AI augmentation in each step and combines the two steps into a single job, resulting in a total cost of $(2+2)(2+2) = 16$ and requiring a worker of skill $4$.
   Note that this is preferable to combining the two steps into a single AI chain, which would result in a total cost of $(2)( (1/4) \times (1/8)^{-1} \times (1/8)^{-1} ) = 32$.
\end{example}



\end{comment}

\subsection{Connection to the Division of Labor Literature}

The job-design problem developed above connects to the literature on the division of labor within a firm.
Applying the Coase--Williamson logic that firm boundaries are shaped by coordination, transaction, and governance costs \citep{coase1937nature, williamson1971vertical, williamson1979transaction} to the ``boundary of the job,'' the firm chooses how many and which tasks to bundle into a worker's role, determining the degree of worker specialization endogenously \citep{murphy1986specialization, becker1992division, dessein2006adaptive}.
The hand-off costs in our model are an intra-firm analogue of transaction costs that AI does not directly reduce, since they capture coordination frictions arising from tacit or social knowledge that is not easily executable by AI \citep{deming2017growing}.\footnote{
\cite{ide2025ai} provides a related view on how AI affects the structure of work, embedding AI as an autonomous problem solver in a knowledge hierarchy model \citep{garicano2000hierarchies, garicano2006organization} and studying how it reshapes spans of control within the firm.
}
Even so, AI chaining restructures the division of labor on two margins: on the intensive margin, by pulling previously manual steps ``under the hood'' of a single AI-executed task; and on the extensive margin, by lowering the cost of linking adjacent steps across a job boundary and thereby making it optimal to redraw that boundary \citep{autor2025expertise}.
